\documentclass[conference]{IEEEtran}

\IEEEoverridecommandlockouts


\usepackage[utf8]{inputenc}
\usepackage[T1]{fontenc}
\usepackage{cite}
\usepackage{amsmath,amssymb,amsthm}
\usepackage{mathtools}
\usepackage{bm}
\usepackage{booktabs}
\usepackage{array}
\usepackage{enumitem}
\usepackage{graphicx}
\usepackage{xcolor}
\usepackage{hyperref}
\usepackage{cleveref}

\hypersetup{hidelinks}


\newcommand{\R}{\mathbb{R}}
\newcommand{\SO}{\mathrm{SO}(3)}
\newcommand{\dd}{\mathrm{d}}
\newcommand{\phys}{f_{\mathrm{phys}}}
\newcommand{\res}{d_{\theta}}
\newcommand{\latent}{F_{\theta}}
\newcommand{\enc}{E_{\theta}}
\newcommand{\dec}{D_{\theta}}
\newcommand{\Integrator}{\mathcal{I}}
\newcommand{\He}{H_{\epsilon}}
\newcommand{\Ep}{E_p}
\newcommand{\Rel}{R_{\epsilon}}


\title{Physics-Prior Latent World Models for Reliable Long-Horizon UAV Prediction under Out-of-Distribution Wind Disturbances}

\author{
\IEEEauthorblockN{CHETHAN}
\IEEEauthorblockA{
Department of Computational \& Data Sciences\\
Indian Institute of Science\\
Email: chethan1@iisc.ac.in}
}

\begin{document}

\maketitle


\begin{abstract} Learned world models provide a mechanism for predicting future states
without repeatedly interacting with the physical environment. However,
recursive prediction can accumulate error over long horizons, and this
problem can become more severe when the operating conditions differ from
those observed during training. Quadrotor UAVs provide a particularly
interesting setting because their nominal rigid-body dynamics are known,
while aerodynamic effects, wind disturbances, actuator mismatch, and
other unmodeled effects remain difficult to describe analytically. This thesis investigates whether incorporating known quadrotor dynamics
as a physical prior into a latent world model can improve the reliability of long-horizon prediction under out-of-distribution (OOD) wind
disturbances. The proposed model decomposes the continuous-time dynamics into a known physics component and a learned residual,
\[
\dot{s}_t =
f_{\mathrm{phys}}(s_t,u_t)
+
d_{\theta}(h_t,s_t,u_t),
\]
where $s_t$ denotes the physical UAV state and $h_t$ denotes a learned
latent memory state. The resulting dynamics are integrated numerically
and fed back into the latent transition to enable open-loop imagination
without future observations.

The primary evaluation compares a purely learned latent dynamics model
with the proposed physics-prior residual model under identical training
data and evaluation conditions. The main evaluation criterion is not
only one-step prediction accuracy but the rate of error growth during
recursive rollouts. A reliable imagination horizon is defined as the
maximum prediction horizon for which the mean position prediction error
remains below a specified tolerance. Experiments will use simulated
quadrotor trajectories with controlled wind distributions, with wind
magnitude extrapolation serving as the primary OOD setting and
directional and transient disturbances serving as secondary stress
tests. The study will also investigate when a physical prior helps and
when model mismatch can reduce its benefit. Visual observations, uncertainty estimation, and model-predictive control are considered extensions rather than requirements of the core thesis.
\end{abstract}

\begin{IEEEkeywords}
UAV, quadrotor, world models, latent dynamics, physics prior,
residual learning, long-horizon prediction, distribution shift,
wind disturbances, reliable imagination.
\end{IEEEkeywords}


\section{Introduction and Motivation}

\subsection{Background}

Unmanned Aerial Vehicles (UAVs), particularly quadrotors, operate
according to well-established nonlinear rigid-body dynamics. A
simplified rigid-body model can be represented by the translational and
rotational equations

\begin{align}
\dot{\mathbf{p}} &= \mathbf{v}, \\
m\dot{\mathbf{v}}
&=
m\mathbf{g}
+
\mathbf{R}\mathbf{f}_{T}
+
\mathbf{f}_{\mathrm{ext}}, \\
\dot{\mathbf{R}}
&=
\mathbf{R}[\boldsymbol{\omega}]_{\times}, \\
\mathbf{J}\dot{\boldsymbol{\omega}}
&=
\boldsymbol{\tau}
-
\boldsymbol{\omega}
\times
(\mathbf{J}\boldsymbol{\omega})
+
\boldsymbol{\tau}_{\mathrm{ext}},
\end{align}

where $\mathbf{p}$ and $\mathbf{v}$ denote position and velocity,
$\mathbf{R}\in SO(3)$ denotes the attitude rotation matrix,
$\boldsymbol{\omega}$ denotes body-frame angular velocity, $m$ is the
vehicle mass, and $\mathbf{J}$ is the body-frame inertia matrix.
The vectors $\mathbf{f}_{T}$ and $\boldsymbol{\tau}$ denote the total
thrust and control torque, respectively. For an idealized quadrotor, the total thrust magnitude can be related
to the rotor angular velocities by

\begin{equation}
f_T = \sum_{i=1}^{4} k_f \Omega_i^2,
\end{equation}


where $k_f$ is the thrust coefficient and $\Omega_i$ is the angular
velocity of rotor $i$. The corresponding body torques are determined
by the quadrotor geometry, rotor configuration, and motor rotation
directions and can be represented through a control-allocation matrix. These equations provide substantial prior knowledge about how the
vehicle evolves. However, the nominal rigid-body model does not
perfectly describe all effects encountered by a real or high-fidelity
simulated UAV. Important sources of mismatch include

\begin{itemize}
    \item aerodynamic drag,
    \item wind and turbulence,
    \item actuator dynamics and delays,
    \item uncertainty in mass and inertia,
    \item aerodynamic coupling,
    \item rotor and ground effects,
    \item modeling and parameter uncertainty.
\end{itemize}

Consequently, the actual system dynamics can be viewed as the
combination of the nominal physics model and the effects that are not
captured by the nominal model. This decomposition motivates hybrid physics--learning approaches in
which known dynamics are handled analytically while a neural model
learns the remaining discrepancy. Importantly, the learned residual in this thesis is not interpreted as
a direct estimate of wind velocity. Instead, it represents the net
effect of wind, aerodynamic mismatch, actuator mismatch, parameter
uncertainty, and other effects not captured by the nominal model.


\subsection{Why Learned World Models?}

A world model learns an internal representation of an environment and a
transition model that can be recursively used to predict future states.
A generic deterministic latent transition can be written as

\begin{equation}
z_{t+1}
=
F_{\theta}(z_t,u_t).
\end{equation}

A stochastic formulation may instead represent

\begin{equation}
p(z_{t+1}\mid z_t,u_t).
\end{equation}

The main advantage of such models is that future trajectories can be
generated internally without physically executing every candidate
action. This is particularly attractive for UAVs because

\begin{itemize}
    \item physical experiments can be expensive,
    \item dangerous conditions are difficult to collect,
    \item long-horizon decision making requires future prediction,
    \item imagined trajectories can be generated much faster than
    physical experiments,
    \item latent state prediction avoids the requirement to generate
    photorealistic future images.
\end{itemize}

However, recursive prediction introduces an important problem. If

\begin{equation}
\hat{s}_{t+1}
=
F_{\theta}(\hat{s}_t,u_t),
\end{equation}

then the prediction at time $t+1$ becomes the input for predicting
$t+2$. Therefore,

\begin{equation}
\hat{s}_{t+1}
\rightarrow
\hat{s}_{t+2}
\rightarrow
\cdots
\rightarrow
\hat{s}_{t+H}.
\end{equation}

where \(H\) is the horizon length. Even a small one-step error can consequently affect all subsequent
predictions. A model that performs well at one-step prediction may
therefore become unreliable at long imagination horizons.


\subsection{Distribution Shift}

The problem becomes more challenging when the conditions during
evaluation differ from those used during training. Let the wind disturbance during training be sampled from

\begin{equation}
\mathbf{w}\sim P_{\mathrm{train}}(\mathbf{w}).
\end{equation}

During evaluation, the model may encounter

\begin{equation}
\mathbf{w}\sim P_{\mathrm{test}}(\mathbf{w}),
\qquad
P_{\mathrm{test}}\neq P_{\mathrm{train}}.
\end{equation}

For UAVs, distribution shifts can arise from

\begin{itemize}
    \item stronger wind magnitudes,
    \item previously unseen wind directions,
    \item different temporal characteristics,
    \item sudden gusts,
    \item moderate changes in vehicle parameters.
\end{itemize}

A purely learned model must extrapolate the complete dynamics under
these conditions. In contrast, some structural relationships, such as
gravity, thrust generation, translational kinematics, and rotational
dynamics, remain valid even when the disturbance distribution changes.

This motivates the central hypothesis of this thesis: if the known
physics explains a sufficiently large portion of the vehicle dynamics,
then the learned model only needs to represent the remaining
discrepancy. This may reduce recursive error growth under distribution
shift.


\subsection{Research Problem}

The central problem addressed by this thesis is:

\begin{quote}
\itshape
Can known quadrotor physical dynamics, incorporated as a prior in a
latent world model, improve the reliability of recursive long-horizon
prediction under out-of-distribution wind disturbances?
\end{quote}

The thesis deliberately focuses on the predictive dynamics problem
rather than requiring successful autonomous navigation or reinforcement
learning.

The central scientific question can therefore be stated more
concisely as

\begin{quote}
\itshape
Does incorporating known UAV physics slow the growth of prediction
error during long-horizon imagination under OOD wind?
\end{quote}


\section{Literature Review}


\subsection{Physics-Based Quadrotor Dynamics}

Quadrotor dynamics have traditionally been modeled using first-principles rigid-body mechanics. Such models provide physically interpretable predictions and are widely used in simulation and control. Their main limitation is that the analytical model is only an approximation of the complete vehicle dynamics. Aerodynamic forces, wind, actuator dynamics, and uncertain parameters can introduce
discrepancies between the model and the actual system. This motivates augmenting analytical models with data-driven components rather than attempting to replace the physical model entirely.


\subsection{Physics-Informed Learning for UAV Dynamics}

Physics-informed learning incorporates known physical structure into
neural models. Gu et al.~\cite{gu2024physics} investigated
physics-informed neural modeling for quadrotor dynamics and incorporated
physical information into the learning process. Their results under
aerodynamic effects and periodic wind demonstrate that physical
knowledge can improve learned UAV dynamics and physical consistency. Such approaches demonstrate the value of physical priors, but a
physics-informed loss and a physics-based transition are not
necessarily equivalent. In a loss-based formulation, physical
knowledge may constrain the solution during training while the learned
transition itself remains a neural approximation. The present thesis
instead investigates the more explicit formulation in which the known
physical dynamics directly participate in the transition used during
imagination.


\subsection{Hybrid Physics--Learning and Residual Dynamics}

A closely related research direction combines a first-principles model
with a learned residual. The general formulation is

\begin{equation}
\dot{s}_t
=
f_{\mathrm{phys}}(s_t,u_t)
+
d_{\theta}(s_t,u_t).
\end{equation}

The physical model explains the known portion of the dynamics, while
the neural component learns the discrepancy. NeuroBEM by Bauersfeld et al.~\cite{bauersfeld2021neurobem} provides a
representative example of this approach for quadrotor modeling. The
method combines a first-principles aerodynamic model with learned
components to capture effects that are difficult to model analytically.
The results demonstrate that hybrid modeling can provide advantages
over both purely analytical and purely learned alternatives. More recent work has also applied physics-guided residual learning to
UAV trajectory prediction. Islam et al.~\cite{islam2026physics}
combine physics-based trajectory prediction with learned residual
corrections and physical feasibility constraints. These studies establish that residual learning is a viable mechanism
for combining analytical dynamics with learned corrections. However,
they do not by themselves answer the question of how such hybrid
dynamics affect the reliability of a recurrent latent world model
during long open-loop imagination.


\subsection{UAV Dynamics under Wind and Distribution Shift}

Wind represents an important source of model uncertainty in UAV
dynamics. It can change both the magnitude and direction of forces
acting on the vehicle and can produce significant differences between
training and deployment conditions. PI-WAN by Wang et al.~\cite{wang2025piwan} investigates physics-informed
wind-adaptive quadrotor dynamics prediction under varying wind,
payload, and external perturbation conditions. The work demonstrates
that incorporating physical information can improve prediction and
closed-loop tracking under previously unseen conditions. This work is particularly relevant because it establishes the value of
physics-informed modeling under environmental distribution shift.
However, its primary objective is wind-adaptive dynamics prediction and
control rather than measuring the reliable imagination horizon of a
latent world model.


\subsection{Latent World Models}

World Models~\cite{ha2018worldmodels} established the idea of learning
compact internal representations of environments together with
predictive dynamics. PlaNet~\cite{hafner2019planet} introduced a
recurrent stochastic latent state-space model for planning, while
Dreamer~\cite{hafner2020dream} demonstrated that behavior can be learned
using trajectories imagined inside a learned latent model. DreamerV3
further demonstrated the scalability of the world-model paradigm
across diverse environments~\cite{hafner2023mastering}. A generic latent world model can be represented as

\begin{equation}
e_t = E_{\theta}(o_t),
\end{equation}

followed by

\begin{equation}
h_{t+1}
=
F_{\theta}(h_t,e_t,u_t).
\end{equation}

The model can subsequently be rolled forward without future
observations. The central limitation relevant to this thesis is recursive error
accumulation. The quality of a long imagined trajectory depends not
only on the one-step transition but also on the stability of repeated
application of the learned transition.


\subsection{World Models for UAVs}

Recent work has increasingly explored learned world models for aerial navigation and UAV prediction. Navigation-oriented models such as AirDreamer~\cite{liu2026airdreamer} and WorldFly~\cite{zheng2026worldfly} world-model approaches demonstrate the potential of learned predictive representations for aerial decision making. More recent work such as SkyJEPA~\cite{skyjepa2026} specifically
investigates long-horizon world modeling for quadrotors. This is
important because it demonstrates that long-horizon learned predictive
models for UAVs are themselves an active research direction. The present thesis therefore does not claim that UAV world modeling is
new. Instead, it studies a narrower question: whether explicit
quadrotor physical structure can improve the reliability of a latent
world model when it is recursively rolled forward under OOD
disturbances.


\subsection{Long-Horizon Reliability}

A key limitation of evaluating learned dynamics using only one-step metrics is that one-step accuracy does not necessarily predict long-horizon reliability ~\cite{zidan2026worldmodels}.

For recursive prediction,

\begin{equation}
\hat{s}_{t+k}
=
F_{\theta}^{(k)}(s_t,u_{t:t+k-1}),
\end{equation}

and the prediction error may increase with $k$. Therefore, a meaningful evaluation should examine prediction error as a function of horizon,

\begin{equation}
E(H)
=
\frac{1}{H}
\sum_{k=1}^{H}
e_{t+k},
\end{equation}

rather than reporting only one-step RMSE. For UAVs, position, velocity, and attitude should be evaluated
separately because they have different physical units and scales.


\subsection{Research Gap}

The existing literature establishes several important facts:

\begin{enumerate}
    \item known physical structure can improve learned UAV dynamics;
    \item residual learning can capture discrepancies between analytical
    and observed dynamics;
    \item physics-informed approaches can improve robustness under wind
    and environmental changes;
    \item latent world models can support recursive imagined rollouts;
    \item UAV world models can support long-horizon predictive reasoning.
\end{enumerate}

However, these directions have largely been evaluated separately.
In particular, there is comparatively limited systematic analysis of
the following controlled comparison:

\begin{itemize}
\item \textbf{Pure learned latent dynamics}
\item \textbf{Physics-prior latent dynamics}
\end{itemize}


when both are evaluated using recursive long-horizon prediction under
controlled OOD wind disturbances. The contribution of this thesis is therefore not the claim that
physics-informed learning, residual dynamics, or latent world models are
individually novel. Instead, the thesis aims to provide a controlled
investigation of their interaction and determine whether the physical
prior changes the rate of long-horizon error growth under distribution
shift.


\section{Research Questions and Hypotheses}


The thesis is organized around one primary research question and a
small number of secondary questions.

\subsection{Primary Research Question}

\begin{quote}
\textbf{RQ1:}
Does incorporating known quadrotor dynamics into a latent world model
improve long-horizon predictive reliability under OOD wind
disturbances?
\end{quote}

\subsection{Secondary Research Questions}

\begin{quote}
\textbf{RQ2:}
Under which OOD wind conditions does the physics prior provide the
largest benefit?
\end{quote}

\begin{quote}
\textbf{RQ3:}
How sensitive is the benefit of the physics prior to the accuracy of
the nominal physical model and the amount of training data?
\end{quote}

\begin{quote}
\textbf{RQ4:}
Can predictive uncertainty, if incorporated as an extension, provide a
useful indicator of the horizon beyond which imagined trajectories
become unreliable?
\end{quote}

Vision is not a primary research question. It is treated as an optional
extension after the core proprioceptive model has been established.


\subsection{Primary Hypothesis}

The primary hypothesis is

\begin{equation}
\boxed{
H_1:\quad
E_p^{\mathrm{phys}}(H)
<
E_p^{\mathrm{pure}}(H)
}
\end{equation}

\noindent
The physics-prior model is expected to exhibit slower multi-step
prediction error growth under OOD wind disturbances.

An important secondary hypothesis is conditional:

\begin{equation}
\boxed{
H_2:\quad
E_p^{\mathrm{phys}}(H;\delta)
\geq
E_p^{\mathrm{pure}}(H;\delta)
}
\end{equation}

\noindent
for sufficiently large $\delta$, where $\delta$ denotes the magnitude
of nominal physics-model mismatch. The benefit of the physics prior is
therefore expected to decrease as the nominal physics model becomes
sufficiently inaccurate. The thesis does not assume that $H_1$ must hold in every condition.
A negative or conditional result is scientifically meaningful if the
experiments identify the conditions under which the physical prior
fails.


\section{Progress Achieved: Proposed Method}


\subsection{Physical State}

The physical state of the UAV is defined as

\begin{equation}
s_t =
\left[
\mathbf{p}_t,\,
\mathbf{v}_t,\,
\mathbf{R}_t,\,
\boldsymbol{\omega}_t
\right],
\end{equation}

where

\begin{align}
\mathbf{p}_t &\in \R^3, &
\mathbf{v}_t &\in \R^3,\\
\mathbf{R}_t &\in \SO, &
\boldsymbol{\omega}_t &\in \R^3.
\end{align}

Quaternion representation may be used for storage and learning, while
the physical dynamics and attitude error are represented using
rotation matrices or an equivalent manifold-aware representation. The control input is

\begin{equation}
u_t =
[\Omega_{1,t},\Omega_{2,t},\Omega_{3,t},\Omega_{4,t}],
\end{equation}

or the equivalent motor/thrust command provided by the simulator.


\subsection{Latent State}

The latent state $h_t$ represents information needed to predict the
future that is not fully captured by the instantaneous physical state. It is updated from sensor observations through

\begin{equation}
h_t =
F_{\theta}
\left(
h_{t-1},
e_t,
u_{t-1}
\right),
\end{equation}

where $e_t$ is an encoded observation. For the core thesis, the observation is restricted to proprioceptive
information:

\begin{equation}
e_t =
E_{\theta}
\left(
a_t^{\mathrm{IMU}},
\boldsymbol{\omega}_t^{\mathrm{IMU}},
u_t
\right).
\end{equation}

This keeps the main problem focused on dynamics learning rather than
visual representation learning.


\subsection{Physical State Decoder}

The latent state is mapped to a physical state estimate through

\begin{equation}
\hat{s}_t =
D_{\theta}(h_t).
\end{equation}

During training, the decoder is supervised using simulator state
measurements. The predicted state is therefore explicitly interpretable in terms of
position, velocity, attitude, and angular velocity.


\subsection{Physics-Prior Transition}

The central model decomposes the continuous-time dynamics into known
physics and learned residual dynamics:

\begin{equation}
\boxed{
\dot{\hat{s}}_t
=
f_{\mathrm{phys}}
(\hat{s}_t,u_t)
+
d_{\theta}
(h_t,\hat{s}_t,u_t).
}
\end{equation}

The first term represents the known nominal quadrotor dynamics. The
second term represents effects that are not adequately captured by the
nominal model. The learned residual captures effects not represented by the nominal physics model:

\begin{equation}
d_{\theta}
\approx
d_{\mathrm{wind}}
+
d_{\mathrm{aero}}
+
d_{\mathrm{param}}
+
d_{\mathrm{other}}.
\end{equation}

\noindent
Here, $d_{\mathrm{wind}}$, $d_{\mathrm{aero}}$, $d_{\mathrm{param}}$,
and $d_{\mathrm{other}}$ denote the contributions of wind effects,
aerodynamic mismatch, parameter mismatch, and other unmodeled effects,
respectively. This avoids requiring the network to identify wind velocity from
available observations.


\subsection{Numerical Integration}

The continuous-time dynamics are converted into a discrete transition
using a numerical integrator, preferably fourth-order Runge--Kutta:

\begin{equation}
\tilde{s}_{t+1}
=
\operatorname{RK4}
\left(
\hat{s}_t,
f_{\mathrm{phys}}(\hat{s}_t,u_t)
+
d_{\theta}(h_t,\hat{s}_t,u_t),
\Delta t
\right).
\end{equation}

The numerical integration step provides a structured transition rather
than asking the neural network to learn the complete state evolution
directly.


\subsection{Latent--Physical Coupling}

After predicting the next physical state, the predicted state is fed
back into the latent dynamics:

\begin{equation}
\boxed{
\hat{h}_{t+1}
=
F_{\theta}
\left(
h_t,
\phi(\tilde{s}_{t+1}),
u_t
\right),
}
\end{equation}

where $\phi(\cdot)$ denotes a state embedding.

The overall architecture is therefore

\begin{equation}
\boxed{
h_t
\rightarrow
\text{physical transition}
\rightarrow
\tilde{s}_{t+1}
\rightarrow
h_{t+1}.
}
\end{equation}

The latent state provides memory, while the physical transition
provides the structured evolution of the UAV.


\subsection{Open-Loop Imagination}

During observation, the latent state can use sensor measurements:

\begin{equation}
h_t =
F_{\theta}
(h_{t-1},e_t,u_{t-1}).
\end{equation}

During imagination, future observations are unavailable. The model must
therefore recursively use its own predictions:

\begin{align}
h_t
&\rightarrow
\tilde{s}_{t+1}
\rightarrow
h_{t+1}
\rightarrow
\tilde{s}_{t+2}
\rightarrow
h_{t+2}
\rightarrow
\cdots \\
&\rightarrow
\tilde{s}_{t+H}.
\end{align}

This open-loop rollout is the primary object evaluated by the thesis.


\section{Future Plans: Training Methodology}


\subsection{Dataset Generation}

The initial dataset will be generated in a physics-based simulator such
as Isaac Sim / Isaac Lab / OmniVerse.

Each trajectory will contain, where available,

\begin{itemize}
    \item position,
    \item linear velocity,
    \item attitude,
    \item angular velocity,
    \item IMU acceleration,
    \item IMU angular velocity,
    \item motor commands or rotor speeds,
    \item wind disturbance,
    \item external disturbance force,
    \item simulation parameters.
\end{itemize}

Wind measurements and disturbance forces are recorded for analysis and
for oracle experiments but are not required as inputs to the core
model.


\subsection{Trajectory Generation}

The dataset will consist of trajectories generated under diverse
control sequences and wind conditions.

The control trajectories should include

\begin{itemize}
    \item hovering,
    \item translational motion,
    \item acceleration and deceleration,
    \item yaw motion,
    \item moderate attitude changes,
    \item combined translational and rotational motion.
\end{itemize}

The dataset should be split by complete trajectories rather than by
individual frames in order to avoid temporal leakage.


\subsection{Wind Distribution}

The primary OOD experiment will use wind-magnitude extrapolation.

For example,

\begin{equation}
|\mathbf{w}|_{\mathrm{train}}
\in [0,4]~\mathrm{m/s},
\end{equation}

while

\begin{equation}
|\mathbf{w}|_{\mathrm{test}}
\in (4,8]~\mathrm{m/s}.
\end{equation}

The exact ranges will be selected after verifying simulator stability
and the operating envelope of the chosen UAV.

Secondary OOD conditions include

\begin{enumerate}
    \item unseen wind directions,
    \item different temporal characteristics,
    \item sudden gust disturbances.
\end{enumerate}

Physics-parameter mismatch is treated primarily as an ablation rather
than as a fourth independent main OOD axis.


\subsection{Physics Self-Consistency Study}

Before training the neural models, the nominal physics model will be
evaluated against simulator ground truth.

\begin{center}
\renewcommand{\arraystretch}{1.15}
\begin{tabular}{c l l}
\toprule
\textbf{Case} & \textbf{Condition} & \textbf{Purpose} \\
\midrule
A & No wind & Baseline physics accuracy \\
B & Training-range wind & Disturbance-induced error \\
C & Parameter perturbation & Physics-model mismatch \\
\bottomrule
\end{tabular}
\end{center}

This experiment is important because the usefulness of a physics prior
depends on the quality of the physical model.


\subsection{Training Stages}

The model will be trained progressively.

\paragraph{Stage 1: Physics-only model.}

Evaluate

\begin{equation}
\dot{s}
=
f_{\mathrm{phys}}(s,u)
\end{equation}

without learning a residual.

\paragraph{Stage 2: Pure neural baseline.}

Train a neural transition model

\begin{equation}
\hat{s}_{t+1}
=
F_{\theta}^{\mathrm{pure}}
(\hat{s}_t,u_t,h_t).
\end{equation}

The model should have approximately comparable parameter capacity to
the proposed residual model.

\paragraph{Stage 3: Physics + residual.}

Train

\begin{equation}
\dot{\hat{s}}
=
f_{\mathrm{phys}}
+
d_{\theta}.
\end{equation}

\paragraph{Stage 4: Physics + residual without latent memory.}

Use

\begin{equation}
d_{\theta}(s_t,u_t)
\end{equation}

instead of

\begin{equation}
d_{\theta}(h_t,s_t,u_t).
\end{equation}

This ablation determines whether recurrent latent memory provides
additional benefit.

\paragraph{Stage 5: Long-horizon training.}

Gradually increase the rollout horizon during training rather than
immediately optimizing extremely long trajectories.

A possible curriculum is

\begin{equation}
H_{\mathrm{train}}
=
1
\rightarrow
5
\rightarrow
10
\rightarrow
20
\rightarrow
50.
\end{equation}

The final horizon will be selected based on computational cost and
initial experimental stability.


\subsection{Prediction Loss}

The physical-state prediction loss is decomposed by state component:

\begin{equation}
\mathcal{L}_{\mathrm{state}}
=
\lambda_p\mathcal{L}_p
+
\lambda_v\mathcal{L}_v
+
\lambda_R\mathcal{L}_R
+
\lambda_\omega\mathcal{L}_\omega.
\end{equation}

Position loss:

\begin{equation}
\mathcal{L}_p
=
\left\|
\hat{\mathbf p}_t-\mathbf p_t
\right\|_2^2.
\end{equation}

Velocity loss:

\begin{equation}
\mathcal{L}_v
=
\left\|
\hat{\mathbf v}_t-\mathbf v_t
\right\|_2^2.
\end{equation}

Attitude loss:

\begin{equation}
\mathcal{L}_R
=
d_{\SO}
(\hat{\mathbf R}_t,\mathbf R_t)^2.
\end{equation}

Angular-velocity loss:

\begin{equation}
\mathcal{L}_\omega
=
\left\|
\hat{\boldsymbol{\omega}}_t
-
\boldsymbol{\omega}_t
\right\|_2^2.
\end{equation}


\subsection{Multi-Step Training Loss}

For an imagination horizon $H$, the multi-step loss is

\begin{equation}
\mathcal{L}_{\mathrm{roll}}
=
\sum_{k=1}^{H}
\gamma^{k-1}
\mathcal{L}_{\mathrm{state}}
(\hat{s}_{t+k},s_{t+k}),
\end{equation}

where $\gamma$ controls the relative importance of later predictions.

The final training objective is

\begin{equation}
\mathcal{L}
=
\mathcal{L}_{\mathrm{state}}
+
\lambda_{\mathrm{roll}}
\mathcal{L}_{\mathrm{roll}}.
\end{equation}

A residual regularization term may be investigated as an ablation:

\begin{equation}
\mathcal{L}_{\mathrm{smooth}}
=
\sum_t
\left\|
d_{t+1}-d_t
\right\|_2^2.
\end{equation}

This term will not be assumed to be universally appropriate because
rapid disturbances can legitimately produce rapidly changing
residuals.


\section{Experimental Baselines}


The main experiment is designed to isolate the contribution of the
physical prior.

\begin{center}
\renewcommand{\arraystretch}{1.2}
\begin{tabular}{c l}
\toprule
\textbf{Model} & \textbf{Description} \\
\midrule
B1 & Physics-only model \\
B2 & Pure neural latent dynamics \\
B3 & Physics + learned residual + latent memory \\
B4 & Physics + residual without latent memory \\
B5 & Residual-only model \\
B6 & Physics + oracle disturbance \\
\bottomrule
\end{tabular}
\end{center}

The central comparison is

\begin{equation}
\boxed{
\mathrm{B2}
\quad\text{vs.}\quad
\mathrm{B3}.
}
\end{equation}

Both models should use the same training trajectories, observation
history, prediction targets, optimizer protocol, and approximately
matched model capacity.

The additional baselines serve different purposes:

\begin{itemize}
    \item B1 measures the performance of the analytical model alone.
    \item B4 isolates the contribution of recurrent latent memory.
    \item B5 determines whether explicitly separating physics from the
    learned component is beneficial.
    \item B6 provides an approximate upper-bound/reference condition when
    the disturbance is available.
\end{itemize}


\section{OOD Evaluation Protocol}


\subsection{Primary OOD Regime: Wind Magnitude}

The primary experiment evaluates extrapolation beyond the training
wind range:

\begin{equation}
P_{\mathrm{train}}
:
|\mathbf w|\in[0,W_{\mathrm{train}}],
\end{equation}

while

\begin{equation}
P_{\mathrm{test}}
:
|\mathbf w|\in
(W_{\mathrm{train}},W_{\mathrm{test}}].
\end{equation}

The model is not retrained on the OOD test distribution.


\subsection{Secondary OOD Regime: Direction}

The second experiment evaluates wind directions that are not present, or are underrepresented, during training. This tests whether the physical structure generalizes to a change in
the direction of the disturbance.


\subsection{Secondary OOD Regime: Temporal Disturbances}

A further stress test changes the temporal structure of the disturbance, such as turbulence frequency or correlation time. The purpose is not to establish universal turbulence robustness but to determine whether the observed advantage of the physical prior persists under a qualitatively different disturbance process.


\subsection{Secondary OOD Regime: Gusts}

Sudden transient wind disturbances are introduced as another stress test. The experiment evaluates whether a model trained primarily on continuous disturbances can maintain reliable predictions when exposed
to short-duration disturbances.


\subsection{Physics-Model Mismatch}

The nominal model parameters are perturbed during evaluation:

\begin{equation}
m' = (1+\delta_m)m,
\end{equation}

\begin{equation}
\mathbf{J}'
=
\mathbf{J}
+
\Delta\mathbf{J}.
\end{equation}

Aerodynamic parameters may also be perturbed. This experiment tests an important limitation: A physics prior is useful only if its structural assumptions remain sufficiently accurate.


\section{Evaluation Metrics}


\subsection{One-Step Prediction}

One-step prediction is evaluated using separate physical metrics.

Position:

\begin{equation}
\mathrm{RMSE}_p
=
\sqrt{
\frac{1}{N}
\sum_{i=1}^{N}
\|
\hat{\mathbf p}_i-\mathbf p_i
\|_2^2
}.
\end{equation}

Velocity:

\begin{equation}
\mathrm{RMSE}_v
=
\sqrt{
\frac{1}{N}
\sum_{i=1}^{N}
\|
\hat{\mathbf v}_i-\mathbf v_i
\|_2^2
}.
\end{equation}

Attitude:

\begin{equation}
e_R
=
d_{\SO}
(\hat{\mathbf R},\mathbf R).
\end{equation}

Angular velocity:

\begin{equation}
\mathrm{RMSE}_{\omega}
=
\sqrt{
\frac{1}{N}
\sum_{i=1}^{N}
\|
\hat{\boldsymbol{\omega}}_i
-
\boldsymbol{\omega}_i
\|_2^2
}.
\end{equation}


\subsection{Long-Horizon Position Error}

Because the primary concern is whether the predicted trajectory remains
physically useful over long horizons, position is used as the primary
reliability quantity.

For a trajectory starting at time $t$, define

\begin{equation}
\boxed{
E_p(H)
=
\frac{1}{H}
\sum_{k=1}^{H}
\left\|
\hat{\mathbf p}_{t+k}
-
\mathbf p_{t+k}
\right\|_2.
}
\end{equation}

Velocity and attitude errors are reported separately:

\begin{equation}
E_v(H)
=
\frac{1}{H}
\sum_{k=1}^{H}
\left\|
\hat{\mathbf v}_{t+k}
-
\mathbf v_{t+k}
\right\|_2,
\end{equation}

and

\begin{equation}
E_R(H)
=
\frac{1}{H}
\sum_{k=1}^{H}
d_{\SO}
(
\hat{\mathbf R}_{t+k},
\mathbf R_{t+k}
).
\end{equation}

This avoids combining quantities with incompatible units into a single
state error.


\subsection{Reliable Imagination Horizon}

For a chosen position-error tolerance $\epsilon$, define

\begin{equation}
H_{\epsilon}
=
\max
\left\{
H:
E_p(k)\leq\epsilon,
\ \forall k\leq H
\right\}.
\end{equation}

This gives a direct interpretation: $H_{\epsilon}$ is the maximum horizon for which the predicted UAV
position remains within the specified average error tolerance. Rather than reporting only one threshold, the experiment will evaluate multiple values of $\epsilon$. For example,

\begin{equation}
\epsilon
\in
\{0.25,0.5,1.0\}\;\mathrm{m},
\end{equation}

subject to the scale and operating range of the selected UAV.


\subsection{Reliability Curve}

A second metric measures the fraction of trajectories that remain
reliable at each horizon.

For trajectory $i$,

\begin{equation}
E_p^{(i)}(H)
=
\frac{1}{H}
\sum_{k=1}^{H}
\left\|
\hat{\mathbf p}^{(i)}_{t+k}
-
\mathbf p^{(i)}_{t+k}
\right\|_2.
\end{equation}

Define

\begin{equation}
\boxed{
R_{\epsilon}(H)
=
\frac{1}{N}
\sum_{i=1}^{N}
\mathbf{1}
\left[
E_p^{(i)}(H)\leq\epsilon
\right].
}
\end{equation}

Thus $R_{\epsilon}(H)$ represents the fraction of trajectories whose average position error remains below $\epsilon$ at horizon $H$. This provides a more robust characterization than a single average
error value.


\section{Ablation Studies}


The following ablations will be performed where computationally
feasible.

\subsection{Physics Ablation}

Compare

\begin{equation}
\text{Pure neural}
\quad\text{vs.}\quad
\text{Physics + residual}.
\end{equation}

This is the primary ablation.


\subsection{Latent-Memory Ablation}

Compare

\begin{equation}
d_{\theta}(s_t,u_t)
\end{equation}

against

\begin{equation}
d_{\theta}(h_t,s_t,u_t).
\end{equation}

This determines whether temporal memory is useful for representing
history-dependent disturbances and model discrepancies.


\subsection{Training-Data Ablation}

Train the models using different fractions of the available training
data:

\begin{equation}
\{25\%,50\%,100\%\}.
\end{equation}

This tests whether the physical prior provides a sample-efficiency
advantage.


\subsection{Physics-Mismatch Ablation}

The nominal physical parameters are gradually perturbed.

The resulting relationship

\begin{equation}
\text{physics mismatch}
\rightarrow
\text{prediction performance}
\end{equation}

will be analyzed to determine when the physical prior becomes harmful.



\section{Optional Vision Extension}


The core thesis does not require visual observations. After the proprioceptive model is established, RGB or depth observations may be added through a visual encoder:

\begin{equation}
z_t^V
=
E_V(I_t).
\end{equation}

The proprioceptive representation is

\begin{equation}
z_t^P
=
E_P
(
a_t^{\mathrm{IMU}},
\boldsymbol{\omega}_t^{\mathrm{IMU}},
u_t
).
\end{equation}

The two representations may be fused:

\begin{equation}
e_t
=
\operatorname{Fuse}
(z_t^V,z_t^P).
\end{equation}

The purpose of this extension is only to determine whether visual information provides predictive information beyond proprioception. It is not necessary for demonstrating the central thesis hypothesis.


\section{Experimental Analysis}


The central experimental analysis will compare prediction-error growth as the imagination horizon increases. For each model, plots will be generated for

\begin{equation}
E_p(H)
\quad\text{versus}\quad
H.
\end{equation}

The primary comparison will be made under

\begin{enumerate}
    \item in-distribution wind,
    \item mild OOD wind,
    \item strong OOD wind,
    \item unseen wind direction,
    \item transient gust conditions.
\end{enumerate}

The expected result, if the primary hypothesis is supported, is

\begin{equation}
E_p^{\mathrm{physics+residual}}(H)
<
E_p^{\mathrm{pure}}(H)
\end{equation}

for sufficiently large $H$ under OOD conditions.

Equivalently,

\begin{equation}
H_{\epsilon}^{\mathrm{physics+residual}}
>
H_{\epsilon}^{\mathrm{pure}}.
\end{equation}

However, the experiments will not assume that this inequality holds universally. The benefit of the physics prior is expected to depend on the magnitude
of the wind distribution shift, the amount and diversity of training data, and the degree of mismatch between the nominal physics model and
the true system dynamics.



\section{Expected Scientific Outcomes}


Three broad outcomes are possible.

\subsection{Physics Provides a Consistent Benefit}

If

\begin{equation}
H_{\epsilon}^{\mathrm{physics+residual}}
>
H_{\epsilon}^{\mathrm{pure}}
\end{equation}

across multiple OOD conditions, the results would provide evidence that
incorporating known UAV physics improves long-horizon predictive
reliability.


\subsection{Physics Provides a Conditional Benefit}

The physical prior may help under moderate distribution shifts or
limited data but provide little advantage under severe mismatch. This would indicate that physical priors are beneficial only within an
appropriate operating regime.


\subsection{Physics Does Not Provide a Benefit}

The physics-prior model may perform similarly to or worse than the
purely learned model.

This could occur if

\begin{itemize}
    \item the nominal physics model is inaccurate,
    \item the learned residual is difficult to optimize,
    \item the distribution shift is too severe,
    \item the latent representation is inadequate,
    \item the purely learned model has sufficient data to learn the
    required dynamics.
\end{itemize}

Such a result would still be scientifically useful because it would
identify limitations of physics priors in learned UAV world models.


\section{Scope and Limitations}


\subsection{Required Components}

The minimum viable thesis consists of

\begin{itemize}
    \item UAV simulation,
    \item nominal quadrotor rigid-body dynamics,
    \item proprioceptive observations,
    \item learned residual dynamics,
    \item recurrent latent state,
    \item numerical physical integration,
    \item open-loop imagination,
    \item controlled OOD wind experiments,
    \item long-horizon prediction evaluation,
    \item reliable imagination horizon analysis.
\end{itemize}


\subsection{Extensions}

The following are extensions if sufficient time remains:

\begin{itemize}
    \item uncertainty estimation,
    \item visual observations,
    \item multimodal sensor fusion,
    \item additional vehicle-parameter shifts,
    \item uncertainty-aware model predictive control.
\end{itemize}


\subsection{Not Required}

The following are explicitly outside the core thesis:

\begin{itemize}
    \item real-world flight testing,
    \item perfect wind estimation,
    \item wind-state identification,
    \item photorealistic future-frame generation,
    \item autonomous navigation,
    \item reinforcement-learning policy training,
    \item full autonomous MPC system,
    \item generalization across arbitrary UAV morphologies.
\end{itemize}

The thesis may use multiple configurations of the same quadrotor family
as an additional robustness experiment, but broad morphology
generalization is not a primary claim.


\section{Feasibility}


The proposed work is designed to be achievable within an approximately eight-month M.Tech thesis timeline.

The primary reasons for feasibility are:

\begin{enumerate}
    \item The dataset can be generated entirely in simulation or use existing datasets.
    \item The required physical state is directly available from the
    simulator.
    \item The core model does not require large-scale image generation.
    \item The physical model reduces the amount of dynamics that the
    neural network must learn.
    \item The main experiments use a small number of controlled OOD
    regimes.
    \item The core thesis does not depend on successful RL or autonomous
    navigation.
\end{enumerate}

The core result should be achievable before the final two months,
leaving time for ablations, uncertainty experiments, analysis, and
thesis writing.


\section{Timeline}


\begin{center}
\renewcommand{\arraystretch}{1.2}
\begin{tabular}{c l c}
\toprule
\textbf{Phase} & \textbf{Work} & \textbf{Duration} \\
\midrule
1 & Literature and simulator setup & 2--3 weeks \\
2 & UAV dynamics and dataset generation & 2--3 weeks \\
3 & Physics calibration and baseline & 2 weeks \\
4 & Pure neural world-model baseline & 2--3 weeks \\
5 & Physics + residual model & 3 weeks \\
6 & Latent memory and open-loop imagination & 3 weeks \\
7 & Primary OOD wind experiments & 3 weeks \\
8 & Ablations and reliability analysis & 3 weeks \\
9 & Optional uncertainty / vision extension & 2--3 weeks \\
10 & Final analysis, writing & 3--4 weeks \\
\bottomrule
\end{tabular}
\end{center}

The minimum viable thesis is complete after Phase~7/8.


\section{Expected Contributions}


The expected contributions of the thesis are:

\begin{enumerate}

    \item \textbf{Physics-prior latent world model:}
    A latent UAV dynamics architecture in which known rigid-body
    dynamics directly participate in the transition and a learned
    residual models unmodeled effects.

    \item \textbf{Controlled long-horizon evaluation:}
    A systematic evaluation of recursive prediction error as a function
    of imagination horizon.

    \item \textbf{OOD wind benchmark:}
    A controlled simulation protocol for evaluating learned UAV
    dynamics under wind conditions outside the training distribution.

    \item \textbf{Reliable imagination horizon:}
    A physically interpretable metric describing how far into the future
    the model can maintain an acceptable level of prediction error.

    \item \textbf{Analysis of when physics helps:}
    An empirical study of the interaction between physical-model
    accuracy, training-data quantity, distribution shift, and
    long-horizon prediction reliability.

\end{enumerate}

The contribution is therefore primarily empirical and methodological:
the thesis seeks to characterize the value and limitations of physical
priors in latent UAV world models rather than claim that the individual
components are themselves new.


\section{Final Thesis Statement}


The proposed thesis investigates whether known physical structure can
improve the reliability of learned latent world models under
distribution shift.

The central model is

\begin{equation}
\boxed{
\dot{\hat{s}}_t
=
f_{\mathrm{phys}}
(\hat{s}_t,u_t)
+
d_{\theta}
(h_t,\hat{s}_t,u_t).
}
\end{equation}

The resulting physical prediction is integrated and fed back into the
latent transition:

\begin{equation}
\boxed{
\hat{h}_{t+1}
=
F_{\theta}
\left(
h_t,
\phi(\tilde{s}_{t+1}),
u_t
\right).
}
\end{equation}

The central evaluation quantity is the position prediction error

\begin{equation}
\boxed{
E_p(H)
=
\frac{1}{H}
\sum_{k=1}^{H}
\left\|
\hat{\mathbf p}_{t+k}
-
\mathbf p_{t+k}
\right\|_2.
}
\end{equation}

The corresponding reliable imagination horizon is

\begin{equation}
\boxed{
H_{\epsilon}
=
\max
\left\{
H:E_p(H)\leq\epsilon
\right\}.
}
\end{equation}

The central research question is therefore:

\begin{quote}
\itshape
How far can a learned UAV world model reliably imagine under wind
disturbances that were not observed during training, and can known
quadrotor physics extend that reliable horizon?
\end{quote}

The thesis does not assume that the physics prior will always help. Instead, it seeks to determine whether, when, and why incorporating known physical structure changes the reliability of long-horizon imagined UAV trajectories.



\begin{thebibliography}{99}

\bibitem{ha2018worldmodels}
D.~Ha and J.~Schmidhuber,
``World Models,''
arXiv preprint arXiv:1803.10122, 2018.

\bibitem{hafner2019planet}
D.~Hafner, T.~Lillicrap, I.~Fischer, R.~Villegas, D.~Ha, H.~Lee,
and J.~Davidson,
``Learning Latent Dynamics for Planning from Pixels,''
in \emph{Proc. International Conference on Machine Learning (ICML)},
vol.~97, pp.~2555--2565, 2019.

\bibitem{hafner2020dream}
D.~Hafner, T.~Lillicrap, J.~Ba, and M.~Norouzi,
``Dream to Control: Learning Behaviors by Latent Imagination,''
in \emph{Proc. International Conference on Learning Representations
(ICLR)}, 2020.

\bibitem{hafner2023mastering}
D.~Hafner, J.~Pasukonis, J.~Ba, and T.~Espeholt,
``Mastering Diverse Domains through World Models,''
arXiv preprint arXiv:2301.04104, 2023.

\bibitem{gu2024physics}
W.~Gu, S.~Primatesta, and A.~Rizzo,
``Physics-informed neural network for quadrotor dynamical modeling,''
\emph{Robotics and Autonomous Systems}, vol.~171, Art.~no.~104569,
2024.

\bibitem{bauersfeld2021neurobem}
L.~Bauersfeld, E.~Kaufmann, P.~Foehn, S.~Sun, and D.~Scaramuzza,
``NeuroBEM: Hybrid Aerodynamic Quadrotor Model,''
in \emph{Proc. Robotics: Science and Systems (RSS)}, 2021,
doi: 10.15607/RSS.2021.XVII.042.

\bibitem{wang2025piwan}
M.~Wang, B.~Wang, Y.~Niu, and C.~Wang,
``PI-WAN: A Physics-Informed Wind-Adaptive Network for Quadrotor
Dynamics Prediction in Unknown Environments,''
in \emph{Proc. IEEE/RSJ International Conference on Intelligent Robots
and Systems (IROS)}, pp.~5413--5419, 2025,
doi: 10.1109/IROS60139.2025.11247234.

\bibitem{islam2026physics}
M.~A.~Islam, S.~Förster, T.~Zhang, \emph{et al.},
``Physics-guided residual learning for phase-aware UAV trajectory
prediction in urban environments,''
\emph{Scientific Reports}, vol.~16, Art.~no.~27073, 2026,
doi: 10.1038/s41598-026-67222-5.

\bibitem{skyjepa2026}
P.~Rao, W.~Zhang, R.~Balestriero, Y.~LeCun, and G.~Loianno,
``SkyJEPA: Learning Long-Horizon World Models for Zero-Shot
Sim-to-Real Control of Quadrotors,''
arXiv preprint arXiv:2606.23444, 2026.

\bibitem{liu2026airdreamer}
Z.~Liu, A.~Yang, C.~Yang, R.~An, C.~Gao, and G.~Zhou,
``AirDreamer: Generalist Drone Navigation with World Models,''
arXiv preprint arXiv:2606.03252, 2026.

\bibitem{zheng2026worldfly}
S.~Zheng, K.~Li, W.~Zhang, Y.~Meng, C.~Gao, X.~Chen, Y.~Li, and
X.-P.~Zhang,
``WorldFly: A World-Model-Based Vision-Language-Action Model for UAV
Navigation,''
arXiv preprint arXiv:2606.06147, 2026.

\bibitem{zidan2026worldmodels}
A.~H.~Zidan, Y.~Pan, H.~Jiang, R.~Yan, W.~Ruan, Z.~Wu, L.~Chen,
W.~You, X.~Li, B.~Chen, H.~Hu, P.~Wang, S.~Liu, J.~Zhang, S.~Li,
Z.~Liu, Y.~Bao, L.~Zhao, L.~Sun, D.~Zhu, X.~Li, J.~Lv, Q.~Li,
W.~Liu, T.~Liu, and W.~Zhang,
``World Models: A Comprehensive Survey of Architectures, Methodologies,
Reasoning Paradigms, and Applications,''
arXiv preprint arXiv:2606.00133, 2026.

\end{thebibliography}

\end{document}